\documentclass[11pt]{article}
\usepackage[margin=1in]{geometry}
\usepackage[T1]{fontenc}
\usepackage{lmodern}
\usepackage{amsmath,booktabs,hyperref,microtype}
\hypersetup{hidelinks}
\makeatletter\let\tableinput\@@input\makeatother
% Generated; do not edit. Sharp k-of-n, interval-robust.
\newcommand{\CatAttackVersion}{17.1}
\newcommand{\CatCoreKFourGuar}{yes}
\newcommand{\CatCoreKOneGuar}{no}
\newcommand{\CatCoreKThreeGuar}{yes}
\newcommand{\CatCoreKTwoGuar}{yes}
\newcommand{\CatEmptyKFourPct}{9.9}
\newcommand{\CatEmptyKOnePct}{44.9}
\newcommand{\CatEpsilon}{5}
\newcommand{\CatMinGuardKFour}{1}
\newcommand{\CatMinGuardKFourCipher}{1}
\newcommand{\CatMinGuardKOne}{3}
\newcommand{\CatMinGuardKOneCipher}{3}
\newcommand{\CatMinGuardKThree}{1}
\newcommand{\CatMinGuardKThreeCipher}{1}
\newcommand{\CatMinGuardKTwo}{2}
\newcommand{\CatMinGuardKTwoCipher}{1}


\title{A sharp distribution-free certificate for catastrophic multi-service outage\\
under security-control portfolios}
\author{}
\date{2026-09-20 (working note, version 0.1.0)}

\begin{document}
\maketitle

\begin{abstract}
A hospital's concern is not that \emph{some} clinical service is disrupted but that
\emph{several critical services fail at once}. We certify exactly that: the
probability that at least $k$ of the $n$ clinical services suffer simultaneous
sustained outage stays below a level $\varepsilon$, under an ATT\&CK-driven attack
model and interval uncertainty. The bound on the $k$-of-$n$ event is the sharp
distribution-free maximum over all joint outcomes consistent with the per-service
marginals -- a small linear program over the $2^n$ outcome atoms -- which replaces
the loose union bound used previously (its $k{=}1$ special case). Because the
catastrophic event is rarer than "any one service", it is far more certifiable: at
$\varepsilon=\CatEpsilon\%$ the minimum guaranteed control count falls from
\CatMinGuardKOne{} controls for $k{=}1$ to \CatMinGuardKThree{} for $k{=}3$, and the
hospital-core bundle, \emph{not} certified for "at least one service"
(\CatCoreKOneGuar), \emph{is} certified for the two-, three-, and all-service
catastrophic events (\CatCoreKTwoGuar/\CatCoreKThreeGuar/\CatCoreKFourGuar). The aggregation
bound is classical (Frechet/Boole/linear programming); the contribution is the
sharp, clinically meaningful certificate, not a new theorem.
\end{abstract}

\section{The k-of-n bound}
Given only the per-service outage marginals $p_1,\dots,p_n$ (here
$p_s = R\,d_s$: kill-chain impact reachability times a service's degradation
probability), the largest possible probability that at least $k$ services fail
together, over every joint distribution with those marginals, is
\[
\max_{\text{joints}} \Pr\!\Big(\textstyle\sum_s \mathbf 1[s\text{ down}] \ge k\Big),
\]
a linear program over the $2^n$ outcome atoms with the marginals and total mass as
equality constraints. Its optimum is attained by an explicit joint, so the bound is
sharp. It equals the union bound $\min(1,\sum_s p_s)$ at $k=1$, equals $\min_s p_s$
at $k=n$, is monotone non-decreasing in each marginal, and is strictly tighter than
the union (and than the Markov bound $\sum p_s/k$) for $k\ge 2$. Monotonicity in the
marginals -- which are themselves monotone in the model parameters -- means the
worst case over the interval uncertainty set is still an interval corner, so the
sharp bound drops into the existing corner certificate unchanged.

\section{Result}
Under the ATT\&CK-driven hospital (v\CatAttackVersion) with the evidence-anchored
intervals, the undefended probability of "at least one" clinical service in
sustained outage is \CatEmptyKOnePct\%, but the all-four catastrophic event is only
\CatEmptyKFourPct\%. The certifiable control burden falls sharply with the
catastrophic threshold (Table~\ref{tab:catsum}): at $\varepsilon=\CatEpsilon\%$,
guaranteeing "at least one service" needs \CatMinGuardKOne{} controls, "at least two"
needs \CatMinGuardKTwo{}, and "at least three" or "all four" need \CatMinGuardKThree{}.
The hospital-core bundle shows the same lesson at the portfolio level
(Table~\ref{tab:catnamed}): it cannot certify the single-service event but does
certify every genuinely catastrophic one. Grounding the degradation in the real
CIPHER corpus lowers the burden further (final column of Table~\ref{tab:catsum}).
The policy reading: a modest, well-chosen portfolio can \emph{provably} bound the
risk of a system-wide clinical collapse even when it cannot bound every isolated
disruption.

\begin{table}[t]\centering
\caption{Minimum guaranteed control count at $\varepsilon=\CatEpsilon\%$ and
undefended worst-case probability, by catastrophic threshold $k$, under assumed and
CIPHER-derived degradation.}\label{tab:catsum}
\begin{tabular}{@{}rrrrr@{}}\toprule
$k$ & Undefended (\%) & Min controls & Undefended, CIPHER (\%) & Min controls, CIPHER\\\midrule
\tableinput table_catastrophic_summary_rows.tex
\bottomrule\end{tabular}
\end{table}

\begin{table}[t]\centering
\caption{Named control bundles: whether each is guaranteed at
$\varepsilon=\CatEpsilon\%$ for the $k$-of-four catastrophic event (assumed
degradation).}\label{tab:catnamed}
\begin{tabular}{@{}lrllll@{}}\toprule
Bundle & Size & $k{=}1$ & $k{=}2$ & $k{=}3$ & $k{=}4$\\\midrule
\tableinput table_catastrophic_named_rows.tex
\bottomrule\end{tabular}
\end{table}

\section{Limitations}
The sharp aggregation bound is classical; this note contributes its use to certify
the clinically defined catastrophic event and its composition with the interval
certificate. The marginals inherit every limitation of the underlying model:
ATT\&CK \texttt{mitigates} edges are qualitative, effectiveness and degradation are
evidence-anchored intervals rather than measurements, and there is no incident
validation. The bound uses only marginals, so it is agnostic to the true
dependence between service failures -- deliberately, since that dependence is
unknown; a validated dependence model could only tighten it.

\paragraph{Reproducibility.} All numbers are generated from committed tables; the
analysis is deterministic and carries a provenance manifest over the pinned ATT\&CK
bundle and the CIPHER corpus.

\end{document}
